\documentclass[11pt,a4paper]{article}
\usepackage[utf8]{inputenc}
\usepackage[T1]{fontenc}
\usepackage[brazil]{babel}
\usepackage[a4paper,margin=2.3cm]{geometry}
\usepackage{enumitem}
\usepackage{hyperref}
\usepackage{xcolor}
\setlength{\parindent}{0pt}
\setlength{\parskip}{5pt}
\sloppy
\setlength{\emergencystretch}{3em}
\hypersetup{colorlinks=true,urlcolor=blue,linkcolor=black}
\newcommand{\cabecalho}[2]{\begin{center}{\small \textbf{Sprint 2 -- #1}}\\[1mm]\textit{Inventário de Ativos e Vulnerabilidades}\\Matheus Moreli Durão Silva -- 12621CBS217\\Cibersegurança -- Bacharelado -- Noturno\end{center}\hrule\vspace{4mm}}
\newcommand{\bloco}[1]{\par\vspace{2mm}\textbf{#1}\par}
\newcommand{\codigo}[1]{\texttt{#1}}

\begin{document}\cabecalho{S2\_11 -- Exclusão de Ativos}{}
\bloco{Política de exclusão} Exclusão é uma operação destrutiva. O sistema identifica claramente o alvo com ID e nome, avisa que suas vulnerabilidades também serão removidas e exige a palavra \codigo{SIM}. Qualquer outra resposta cancela a operação.
\bloco{Implementação} A tela mostra ``Confirme o alvo: ID 10 - SRV-ARQUIVOS''. Ela converte a confirmação em \codigo{True} ou \codigo{False}; a regra \codigo{excluir\_ativo} só executa \codigo{del ativos[identificador]} quando recebe \codigo{True}. Dessa forma a regra também pode ser testada sem simular teclado.
\bloco{Casos importantes} Se o registro não existe, nada é excluído e a mensagem informa isso. Se o usuário cancela, a função devolve \codigo{False} e o dicionário preserva a chave. Se confirma, a função devolve \codigo{True} e a chave deixa de existir.
\bloco{Como testar} O teste de exclusão tenta primeiro \codigo{False} e depois \codigo{True}, verificando a presença e a ausência da chave, respectivamente.
\bloco{Checklist atendido} Alvo identificado, confirmação obrigatória, cancelamento sem mudança, inexistência tratada e política documentada.
\end{document}
